\section{Preliminaries}\label{sec-prelims}

\headingg{Notation.} Strings are identified with vectors over $\bits$, so that $|Z|$ denotes the length of a string $Z$ and $Z_i$ denotes its $i$-th bit. By $x\|y$ we denote the concatenation of strings $x,y$. If $S$ is a finite set, then $|S|$ denotes its size and we let $x \getsr S$ denote picking an element of~$S$ uniformly at random and assigning it to $x$. We use $[n]$ to denote the set $\{1, \dots, n\} \subset \N$. By $\secp \in \N$ we denote the security parameter and by $\usecp$ its unary encoding. A function $\nu \colon \N \to \R$ is negligible if for every polynomial $p$ we have $\nu(\secp) \le 1/p(\secp)$ for all sufficiently large $\secp$.

For a function $f \colon A \to B$, we let $\Image(f) := \{f(a) : a \in A\}$, and we let $\mathsf{Size}_f(b) := |f^{-1}(b)|$ for all $b \in B$.

Algorithms may be randomized unless otherwise indicated. If $\advA$ is an algorithm, we let $y\gets \advA^{\procfont{O}_1,\ldots}(x_1,\ldots ; \omega)$ denote running $\advA$ on inputs $x_1,\ldots$ and coins $\omega$, with oracle access to $\procfont{O}_1,\ldots$, and assigning the output to $y$. Moreover, by $y\getsr \advA^{\procfont{O}_1,\ldots}(x_1,\ldots)$ we denote picking $\omega$ uniformly at random and letting $y\gets \advA^{\procfont{O}_1,\ldots}(x_1,\ldots ; \omega)$. We let $\OutSet{\advA^{\procfont{O}_1,\ldots}(x_1,\ldots)}$ denote the set of all possible outputs of $\advA$ when run on inputs $x_1,\ldots$ and with oracle access to $\procfont{O}_1,\ldots$. Running time is worst-case. We use $\bot$ as a special symbol to denote rejection, and it is assumed not to be in $\bits^*$.

\heading{Games.} We use the code-based game-playing framework of BR~\cite{EC:BelRog06}. By $\Pr[\Gm\Rightarrow y]$ we denote the probability that the execution of game $\Gm$ results in output $y$. In games, integer variables, set variables, boolean variables and string variables are assumed initialized, respectively, to $0$, the empty set $\emptyset$, the boolean $\false$ and $\bot$.

\subsection{Signature Schemes and Schnorr Signatures} \label{sec-schnorr}

\heading{Signature schemes.}
A \emph{signature scheme} with message space $\MsgSp$ is a tuple of algorithms $\SigScheme = (\SigKeys, \SigSign, \SigVf)$ that work as follows. The key-generation algorithm $\SigKeys$ outputs a key pair $(\vkey,\skey)$. (We suppress the security parameter for simplicity.) The signing algorithm $\SigSign$ on inputs a signing key $\skey$ and a message $m \in \MsgSp$ outputs a signature $\sig$. The deterministic verification algorithm $\SigVf$ on inputs a verification key $\vkey$, a message $m \in \MsgSp$, and a signature $\sig$ outputs a bit. For correctness, we require that
$$
 \Pr\big[ \SigVf\big(\vkey, m, \SigSign(\skey, m)\big) \Rightarrow 1 \big]
 = 1
$$
for all $(\vkey,\skey) \in \OutSet{\SigKeys}$ and all $m \in \MsgSp$, where the probability is over the coins for $\SigSign$.

\heading{Keyed hash families.} A \emph{keyed hash family} for a group $\GG$ of prime order $p$ and message space $\MsgSp$ is a pair $\HF = (\HFKg, \HFEv)$ that works as follows. The key-generation algorithm $\HFKg$ on input a group element $X \in \GG$ outputs a hash key $\hk$. The deterministic evaluation algorithm $\HFEv$ on inputs $\hk$ and a pair $(R,m) \in \GG \times \MsgSp$ outputs an element of $\Zp$.

\paragraph*{Discussion.} The hash key is allowed to depend on the group element $X$, which in our schemes is the Schnorr verification key. The hash key is generated honestly at key-generation time and placed in the verification key. Our key-recovery construction (Section~\ref{sec-lf-kr}) ignores $X$. The two unforgeability constructions (Sections~\ref{sec-io-uf} and~\ref{sec-adaptive}) embed $X$ in an obfuscated program. Generating the hash key at key-generation time is the same modeling choice as in the full-domain-hash instantiation of Hohenberger, Sahai, and Waters~\cite{EC:HohSahWat14}. It means each user has their own hash function, which is public and part of the verification key.

\heading{Schnorr signatures relative to a hash family.} Let $\GG$ be a cyclic group of prime order $p = |\GG|$, generated by $g$. Let $\HF$ be a keyed hash family for $\GG$ and message space $\MsgSp$. The \emph{Schnorr signature scheme}~\cite{C:FiaSha86,C:Schnorr89} $\SchScheme[\GG,\HF] = (\SchKeys,\SchSign,\SchVf)$ with message space $\MsgSp$ works as follows. Algorithm $\SchKeys$ chooses $x \getsr \Zp$, sets $X \gets g^x$, computes $\hk \getsr \HFKg(X)$, and returns $(\vkey=(X,\hk),\skey=(x,\hk))$. Algorithm $\SchSign$ on input $\skey,m$ chooses $r \getsr \Zp$, sets $R \gets g^r$ and $c \gets \HFEv(\hk, (R, m))$, then returns $\big(R,(r+ cx)\bmod p\big)$. Algorithm $\SchVf$ on inputs $(X,\hk),m,(R,s)$ returns $\big(g^s = R \cdot X^c\big)$ where $c \gets \HFEv(\hk, (R, m))$.
Correctness follows by substitution. We take $\MsgSp = \bits^{\msglen}$ for a polynomial $\msglen = \msglen(\secp)$.

\paragraph*{Discussion.} Restricting to fixed-length messages loses no generality. To sign messages in $\bits^*$, one first hashes the message to $\bits^{\msglen}$ with a collision-resistant hash function whose key is placed in the verification key. For unforgeability, a forgery on a fresh message whose digest collides with a queried message's digest gives a collision. For key recovery, digest collisions do not affect the winning condition at all.

\begin{figure}[t]
\threeCols{0.33}{0.36}{0.27}{
\ExperimentHeader{Game $\gameEUFCMA{\SigScheme}$}

\begin{oracle}{$\Initialize$}
\item $(\vkey,\skey) \getsr \SigKeys$
\item $S \gets \emptyset$
\item Return $\vkey$
\end{oracle}
\ExptSepSpace
\begin{oracleC}{$\SignOracle(m)$}{$m \in \MsgSp$}
\item $\sigma \getsr \SigSign(\skey,m)$
\item $S \gets S \cup \{m\}$
\item Return $\sigma$
\end{oracleC}
\ExptSepSpace

\begin{oracle}{$\Finalize{(m,\sigma)}$}
\item If $m\in S$ return 0
\item Return $\SigVf(\vkey,m,\sigma)$
\end{oracle}
}
{
\ExperimentHeader{Game $\gameEUFSTATIC{\SigScheme}$}

\begin{oracle}{$\Initialize$}
\item $((m_1,\ldots,m_q),m^*,\state) \getsr \advA_0$
\item If $m^* \in \{m_1,\ldots,m_q\}$:
\item \quad $\mathsf{bad} \gets \true$
\item $(\vkey,\skey) \getsr \SigKeys$
\item For $i \in [q]$: $\sigma_i \getsr \SigSign(\skey,m_i)$
\item Return $(\vkey,(\sigma_i)_{i \in [q]},\state)$
\end{oracle}
\ExptSepSpace

\begin{oracle}{$\Finalize{(\sigma)}$}
\item If $\mathsf{bad}$ return 0
\item Return $\SigVf(\vkey,m^*,\sigma)$
\end{oracle}
}
{
\ExperimentHeader{Game $\gameBKCMA{\SigScheme}$}

\begin{oracle}{$\Initialize$}
\item $(\vkey,\skey) \getsr \SigKeys$
\item Return $\vkey$
\end{oracle}
\ExptSepSpace
\begin{oracleC}{$\SignOracle(m)$}{$m \in \MsgSp$}
\item $\sigma \getsr \SigSign(\skey,m)$
\item Return $\sigma$
\end{oracleC}
\ExptSepSpace

\begin{oracle}{$\Finalize{(\skey')}$}
\item Return $(\skey' = \skey)$
\end{oracle}
}
\vspace{-12pt}
	\caption{Games defining EUF-CMA (left), static-message unforgeability EUF-static (middle), and key recovery BK-CMA (right) of a signature scheme $\SigScheme$. In the middle game, the adversary is a pair $\advA = (\advA_0,\advA_1)$, where $\advA_1$ is run on the output of $\Initialize$ and its output is passed to $\Finalize$. In the right game, for Schnorr signatures we identify $\skey$ with $x$, so $\Finalize(x')$ returns $(x' = x)$.}\label{fig-eufcma}
	\vspace{-3pt}\hrulefill
	\vspace{-1ex}
\end{figure}

\begin{samepage}
\heading{Unforgeability and key recovery.} Let $\SigScheme = (\SigKeys, \SigSign,\SigVf)$ be a signature scheme with message space $\MsgSp$. We define three security games, given in Figure~\ref{fig-eufcma}. For an adversary $\advA$, its EUF-CMA advantage is
$$
\AdvEUFCMA{\SigScheme}{\advA}
=\Pr\big[\gameEUFCMA{\SigScheme}(\advA)\Rightarrow1\big].
$$
For $\advA=(\advA_0,\advA_1)$, its \emph{static-message} unforgeability advantage is
$$
\AdvEUFSTATIC{\SigScheme}{\advA}
=\Pr\big[\gameEUFSTATIC{\SigScheme}(\advA)\Rightarrow1\big].
$$
Finally, the \emph{key-recovery} advantage of $\advA$ is
$$
\AdvBKCMA{\SigScheme}{\advA}
=\Pr\big[\gameBKCMA{\SigScheme}(\advA)\Rightarrow1\big].
$$
\end{samepage}

\paragraph*{Discussion.} In EUF-static, the adversary declares its signing messages $m_1,\ldots,m_q$ and its forgery message $m^*$ before seeing the verification key, and then receives all requested signatures at once. This is weaker than the selective-target notion of Hohenberger, Sahai, and Waters~\cite{EC:HohSahWat14}, where only the forgery message is declared before key generation and signing queries remain adaptive. In the key-recovery game BK-CMA, the adversary has an unrestricted, adaptive signing oracle but must output the signing key to win. The name BK-CMA (``breaking the key'') follows Paillier and Vergnaud~\cite{AC:PaiVer05}. Note that BK-CMA and EUF-static are incomparable. The first has a stronger oracle but a harder winning condition.

\subsection{Discrete-Logarithm and Circular Discrete-Logarithm Problems} \label{sec-dl-cdl}

We recall the \emph{discrete-logarithm} (DL) problem.
Let $\GG$ be a group of prime order $p = |\GG|$, generated by $g$. For an adversary $\advA$, we let its DL-advantage against $\GG,g$ be $\AdvDL{\GG,g}{\advA} = \Pr\big[\gameDL{\GG,g}(\advA) \Rightarrow 1\big]$, where the game is in Figure~\ref{fig-cdl}.

For a function $T = T(\secp)$, we say DL is \emph{$T$-hard} in $\GG$ if every probabilistic algorithm running in time at most $T(\secp) \cdot \mathrm{poly}(\secp)$ has negligible DL-advantage. Hardness against PPT algorithms is the special case of constant $T$.

We recall the \emph{circular discrete-logarithm} (CDL) problem, introduced by Cho, Fuchsbauer, O'Neill, and Sefranek~\cite{C:CFOS25}.
Let $\GG$ be as above and let $f \colon \GG \to \Zp$ be an efficient function, called the \emph{conversion function}. For an adversary $\advA$ we let its CDL-advantage against $\GG,g,f$ be $\AdvCDL{\GG,g,f}{\advA} = \Pr\big[\gameCDL{\GG,g,f}(\advA) \Rightarrow 1\big]$, where the game is in Figure~\ref{fig-cdl}.

\begin{figure}[t]
\twoCols{0.4}{0.5}{
\ExperimentHeader{Game $\gameDL{\GG, g}$}

\begin{oracle}{$\Initialize$}
\item $x \getsr \Zp$
\item $h \gets g^x$
\item Return $h$
\end{oracle}
\ExptSepSpace
\begin{oracle}{$\Finalize(x')$}
\item Return $(x = x')$
\end{oracle}
}
{
\ExperimentHeader{Game $\gameCDL{\GG, g, f}$}

\begin{oracle}{$\Initialize$}
\item $x \getsr \Zp$
\item $h \gets g^x$
\item Return $h$
\end{oracle}
\ExptSepSpace
\begin{oracle}{$\Finalize(R, z)$}
\item Return $\big(f(R)\ne 0\:\land\:g^z=R\cdot h^{f(R)}\big)$
\end{oracle}
}
\vspace{-12pt}
	\caption{Games defining DL and circular DL problems.}\label{fig-cdl}
	\vspace{-3pt}\hrulefill
	\vspace{-1ex}
\end{figure}

\paragraph*{Discussion.} The ``circularity'' in CDL is that in the solution equation $R$ occurs both in the base and in the exponent, mirroring the verification equation of Schnorr signatures. CDL is non-interactive and, for a fixed $f$, falsifiable in the sense of Naor~\cite{C:Naor03}. It cannot hold if some nonzero value has a large preimage under $f$; \cite{C:CFOS25} show that in the elliptic-curve generic group model of Groth and Shoup~\cite{EC:GroSho22} this condition on $f$ is also sufficient, and that the ECDSA conversion function $f_{\mathsf{ecdsa}} \colon (x,y) \mapsto x \bmod p$ satisfies it. CDL is random self-reducible~\cite{C:CFOS25}. For our results we require, as in~\cite{C:CFOS25}, that $\mathsf{Size}_f(0)/p$ is negligible and that $f$ is efficient.

\subsection{Tunable Lossy Functions} \label{sec-lf}

We define a lossy function whose image bound is an input to key generation. This lets a reduction choose the bound that gives the best tradeoff between indistinguishability and enumeration time. The definition generalizes the enumerable form of Zhandry's extremely lossy functions~\cite{C:Zhandry16}.

Let $M=2^{\ellin}$, where $\secp\le\ellin=\poly(\secp)$, and let $\Ygrave$ be a range whose elements have polynomial-size descriptions. The lower bound on $\ellin$ is without loss of generality, since an encoding can be padded.

\begin{definition} \label{def-lf}
A \emph{tunable lossy function with enumerable image} is a tuple $\TLF=(\TLFKg,\TLFEv,\TLFEnum)$. Algorithm $\TLFKg$ on input $(\usecp,r,\mathsf{mode})$, where $1\le r\le M$ and $\mathsf{mode}\in\{\mathsf{inj},\mathsf{lossy}\}$, outputs $(K,\mathsl{td})$ in time polynomial in $\secp$. Algorithm $\TLFEv$ on input $K$ and $x\in\bits^{\ellin}$ outputs an element of $\Ygrave$ in time polynomial in $\secp$. Algorithm $\TLFEnum$ on input $\mathsl{td}$ outputs a list of at most $r$ elements of $\Ygrave$; the enumeration information $\mathsl{td}$ includes $r$. There are functions $\nu_{\mathsf{inj}}$ and $\nu_{\mathsf{enum}}$ such that:
\begin{enumerate}
\item \emph{Injectivity:} for every $r\in[M]$, if $(K,\mathsl{td})\getsr\TLFKg(\usecp,r,\mathsf{inj})$, then $\TLFEv(K,\cdot)$ is injective except with probability $\nu_{\mathsf{inj}}(\secp,r)$.
\item \emph{Lossiness and enumeration:} for every $r\in[M]$, if $(K,\mathsl{td})\getsr\TLFKg(\usecp,r,\mathsf{lossy})$ and $S\getsr\TLFEnum(\mathsl{td})$, then
$$
\Image(\TLFEv(K,\cdot))\subseteq S
$$
except with probability $\nu_{\mathsf{enum}}(\secp,r)$. We let $\tau(\secp,r)$ denote the running time of $\TLFEnum$.
\end{enumerate}
For $\mathsf{mode}\in\{\mathsf{inj},\mathsf{lossy}\}$, define
$$
\Pr_{\mathsf{mode}}[\advA;r]
:=
\Pr\bigl[\advA(K)=1:(K,\mathsl{td})\getsr
\TLFKg(\usecp,r,\mathsf{mode})\bigr].
$$
The distinguishing advantage is
$$
\AdvTLF{\TLF,r}{\advA}
:=
\left|\Pr_{\mathsf{inj}}[\advA;r]
-\Pr_{\mathsf{lossy}}[\advA;r]\right|.
$$
\end{definition}

Definition~\ref{def-lf} deliberately leaves the relation among $r$, $\AdvTLF{\TLF,r}{\advA}$, and $\tau(\secp,r)$ explicit. A fixed-parameter lossy function is obtained by choosing one function $r=r(\secp)$.

We say that the injective mode is \emph{universal} if, for every $r,r'\in[M]$, the marginal distributions of $K$ in
$$
(K,\mathsl{td})\getsr\TLFKg(\usecp,r,\mathsf{inj})
\qquad\text{and}\qquad
(K,\mathsl{td})\getsr\TLFKg(\usecp,r',\mathsf{inj})
$$
are identical.

A \emph{strongly efficiently enumerable ELF} is the special case with a universal injective mode in which, for every polynomial time bound $t$ and inverse polynomial $\epsilon$, there is a polynomial $q$ such that every adversary of time at most $t(\secp)$ has $\AdvTLF{\TLF,r}{\advA}\le\epsilon(\secp)$ for every $r\ge q(\secp)$. For polynomial $r$, the functions $\nu_{\mathsf{inj}}(\secp,r)$ and $\nu_{\mathsf{enum}}(\secp,r)$ are negligible, and $\tau(\secp,r)$ is polynomial in $\secp$ and $r$.

For an ELF with domain size $M$, the injective generator runs $\ELFKg(M,M)$ and the lossy generator with bound $r$ runs $\ELFKg(M,r)$.

Zhandry constructs such an ELF from exponential DDH, or more generally from a constant-$k$ exponential $k$-linear assumption~\cite{C:Zhandry16}. The construction first builds a bounded-adversary lossy function and then iterates it at several auxiliary security levels. The enumeration and security properties above follow from Theorem~3.9, Lemma~3.6, and Corollary~3.15 of the full version~\cite{EPRINT:Zhandry16}.

We also fix an efficient deterministic encoding $\enc{\cdot} \colon \GG \times \bits^{\msglen} \to \bits^{\ellin}$ that is injective.

\subsection{Puncturable Pseudorandom Functions} \label{sec-pprf}

\begin{definition} \label{def-pprf}
A \emph{puncturable pseudorandom function} with domain $\bits^{\ellin'}$ and range $\mathcal{Z}$ is a tuple $\PRF = (\PRFKg, \PRFEv, \PRFPunc)$ that works as follows. The key-generation algorithm $\PRFKg$ on input $\usecp$ outputs a key $k$. The deterministic evaluation algorithm $\PRFEv$ on inputs $k$ and $x \in \bits^{\ellin'}$ outputs an element of $\mathcal{Z}$. The puncturing algorithm $\PRFPunc$ on inputs $k$ and a set $T \subseteq \bits^{\ellin'}$ of polynomial size outputs a punctured key $k\{T\}$. We require:
\begin{enumerate}
\item \emph{Functionality preservation:} for all $k \in \OutSet{\PRFKg(\usecp)}$, all polynomial-size $T$, all $k\{T\} \in \OutSet{\PRFPunc(k,T)}$, and all $x \notin T$, it holds that $\PRFEv(k\{T\},x) = \PRFEv(k,x)$.
\item \emph{Pseudorandomness at punctured points:} let $\mathcal{S}$ be any PPT sampler that on input $\usecp$ outputs a polynomial-size set $T$ and state $\state$, independently of the PRF key. For every PPT adversary $\advA$, the advantage
$$
\left|
\Pr\!\left[\advA(\state,T,k\{T\},(\PRFEv(k,x))_{x\in T})\Rightarrow1\right]
-
\Pr\!\left[\advA(\state,T,k\{T\},(u_x)_{x\in T})\Rightarrow1\right]
\right|
$$
is negligible, where $(T,\state)\getsr\mathcal{S}(\usecp)$, $k\getsr\PRFKg(\usecp)$, $k\{T\}\getsr\PRFPunc(k,T)$, and $u_x\getsr\mathcal{Z}$ independently for each $x\in T$.
\end{enumerate}
\end{definition}

Puncturable PRFs follow from one-way functions via the GGM construction~\cite{AC:BonWat13,EC:BoyGilIsh15,CCS:KPTZ13}.

\subsection{Indistinguishability Obfuscation} \label{sec-io}

\begin{definition} \label{def-io}
A PPT algorithm $\iO$ is an \emph{indistinguishability obfuscator} for a circuit class $\{\mathcal{C}_\secp\}$ if the following hold.
\begin{enumerate}
\item \emph{Correctness:} for all $\secp$, all $\CircC \in \mathcal{C}_\secp$, all $\tilde{\CircC} \in \OutSet{\iO(\usecp, \CircC)}$, and all inputs $w$, it holds that $\tilde{\CircC}(w) = \CircC(w)$.
\item \emph{Indistinguishability:} let $\mathcal{S}$ be any PPT sampler that on input $\usecp$ outputs circuits $\CircC_0,\CircC_1\in\mathcal{C}_\secp$ and state $\state$, where the circuits have the same size and compute the same function. For every PPT adversary $\advA$, the advantage
$$
\left|
\Pr[\advA(\state,\iO(\usecp,\CircC_0))\Rightarrow1]
-
\Pr[\advA(\state,\iO(\usecp,\CircC_1))\Rightarrow1]
\right|
$$
is negligible, where $(\CircC_0,\CircC_1,\state)\getsr\mathcal{S}(\usecp)$.
\end{enumerate}
\end{definition}

Existing constructions derive indistinguishability obfuscation from explicit falsifiable assumptions~\cite{STOC:JaiLinSah21,EC:JaiLinSah22}. When we obfuscate a circuit, we pad it to the maximum size of the circuits appearing in the security proof; we do not mention this padding again.

\heading{Subexponential security.} We say a puncturable PRF, an indistinguishability obfuscator, or a problem such as DL or CDL is \emph{subexponentially secure with parameter $\delta$}, for a constant $0 < \delta < 1$, if every possibly non-uniform adversary running in time at most $2^{\secp^{\delta}}$ has advantage at most $2^{-\secp^{\delta}}$ in the corresponding game, for all sufficiently large $\secp$. Theorem~\ref{thm-adaptive} is stated with an explicit advantage bound $\eps$; subexponential security is one way to satisfy it, and quasi-polynomial security suffices when the lossy image bound of the pre-hash is quasi-polynomial (Corollary~\ref{cor-adaptive-qp}). The static-message result uses ordinary PPT security at the level of CDL, iO, and puncturable PRFs. Corollary~\ref{cor-bk} separately uses exponential DDH for Zhandry's ELF construction.
